\documentclass[11pt]{article}
% Shared preamble of the RMC kernel write-ups (% Shared preamble of the RMC kernel write-ups (% Shared preamble of the RMC kernel write-ups (\input{preamble} after \documentclass).
\usepackage[margin=1in]{geometry}
\usepackage{amsmath,amssymb,bm}
\usepackage{graphicx}
\usepackage{booktabs}
\usepackage{hyperref}

\newcommand{\dd}{\,\mathrm{d}}
\newcommand{\Ein}{E}
\newcommand{\Eout}{E'}
\newcommand{\Ecm}{E_{c}}
\newcommand{\mucm}{\mu_{c}}
\newcommand{\mulab}{\mu_0}

\newcommand{\acespec}[1]{J.~L.~Conlin and P.~Romano,
\href{https://mcnp-green.lanl.gov/pdf_files/TechReport_2019_LANL_LA-UR-19-29016_ConlinRomano.pdf}
{\emph{A Compact ENDF (ACE) Format Specification}}, LA-UR-19-29016 (2019),
Sec.~#1}
\newcommand{\ENDFspec}[1]{\href{https://www.nndc.bnl.gov/endfdocs/ENDF-102-2023.pdf}
{\emph{ENDF-6 Formats Manual}} (2023), #1}

% Common closing section: how the verification figures are produced
\newcommand{\verificationmethod}{%
Each figure is produced by \texttt{verification/verify\_laws.py}. For an incident
energy $\Ein$ along $+z$, $10^6$ collisions are sampled with MC/DC's own reaction
sampler (\texttt{sample\_elastic\_scattering}, \texttt{sample\_inelastic\_scattering}
or \texttt{sample\_fission}), and every emitted neutron is histogrammed in
$48 \times 24$ bins of $(\Eout, \mulab)$, with $\mulab$ the lab cosine to $+z$. The
inverted kernel used by RMC is integrated over the same bins (Gauss--Legendre panels;
two-body lines are split where $\mulab^*(\Eout)$ crosses a cosine edge), multiplied by
the number of collisions, and compared with a Pearson-type $\chi^2$ diagnostic
over bins with at least 5 expected counts (the rest pooled into one bin).
The existing tables list the script's $\chi^2/\mathrm{dof}$ and nominal $p$-value,
the emitted neutrons per collision (sampled vs.\ the integral of the
inverted kernel), and the largest relative change of a bin integral when the
quadrature panels are doubled. Left panels: outgoing-energy marginal; middle:
lab-cosine marginal; right: per-bin $z = (\text{sampled} - \text{inverted}) /
\sqrt{\text{inverted}}$.

\noindent These figures check agreement with the MC/DC sampler, not independent
physical validity. Their existing $p$-values are not calibrated goodness-of-fit
probabilities: the script uses $B$ degrees of freedom, whereas a fixed-total
multinomial histogram with $B$ categories and no fitted parameters uses $B-1$.
Random multiplicities and spectrum assignments can also change the collision-level
covariance from independent Poisson counts; $z$ is then only a count diagnostic.
See \href{https://docs.scipy.org/doc/scipy/reference/generated/scipy.stats.chisquare.html}
{the Pearson test's sampling assumptions}.
Energy bins use the sample extrema, so unsampled tails are excluded and the
reported inverted yield is an integral over that observed window. A formal test
needs independently specified bins, explicit tail categories, and a covariance or
resampling model appropriate to independent collisions. Zero expected density
with observed emissions must be reported rather than discarded. Doubling quadrature
panels gives a convergence indicator, not a rigorous error bound. The stored
figures and tables have not been regenerated for these qualifications.}
 after \documentclass).
\usepackage[margin=1in]{geometry}
\usepackage{amsmath,amssymb,bm}
\usepackage{graphicx}
\usepackage{booktabs}
\usepackage{hyperref}

\newcommand{\dd}{\,\mathrm{d}}
\newcommand{\Ein}{E}
\newcommand{\Eout}{E'}
\newcommand{\Ecm}{E_{c}}
\newcommand{\mucm}{\mu_{c}}
\newcommand{\mulab}{\mu_0}

\newcommand{\acespec}[1]{J.~L.~Conlin and P.~Romano,
\href{https://mcnp-green.lanl.gov/pdf_files/TechReport_2019_LANL_LA-UR-19-29016_ConlinRomano.pdf}
{\emph{A Compact ENDF (ACE) Format Specification}}, LA-UR-19-29016 (2019),
Sec.~#1}
\newcommand{\ENDFspec}[1]{\href{https://www.nndc.bnl.gov/endfdocs/ENDF-102-2023.pdf}
{\emph{ENDF-6 Formats Manual}} (2023), #1}

% Common closing section: how the verification figures are produced
\newcommand{\verificationmethod}{%
Each figure is produced by \texttt{verification/verify\_laws.py}. For an incident
energy $\Ein$ along $+z$, $10^6$ collisions are sampled with MC/DC's own reaction
sampler (\texttt{sample\_elastic\_scattering}, \texttt{sample\_inelastic\_scattering}
or \texttt{sample\_fission}), and every emitted neutron is histogrammed in
$48 \times 24$ bins of $(\Eout, \mulab)$, with $\mulab$ the lab cosine to $+z$. The
inverted kernel used by RMC is integrated over the same bins (Gauss--Legendre panels;
two-body lines are split where $\mulab^*(\Eout)$ crosses a cosine edge), multiplied by
the number of collisions, and compared with a Pearson-type $\chi^2$ diagnostic
over bins with at least 5 expected counts (the rest pooled into one bin).
The existing tables list the script's $\chi^2/\mathrm{dof}$ and nominal $p$-value,
the emitted neutrons per collision (sampled vs.\ the integral of the
inverted kernel), and the largest relative change of a bin integral when the
quadrature panels are doubled. Left panels: outgoing-energy marginal; middle:
lab-cosine marginal; right: per-bin $z = (\text{sampled} - \text{inverted}) /
\sqrt{\text{inverted}}$.

\noindent These figures check agreement with the MC/DC sampler, not independent
physical validity. Their existing $p$-values are not calibrated goodness-of-fit
probabilities: the script uses $B$ degrees of freedom, whereas a fixed-total
multinomial histogram with $B$ categories and no fitted parameters uses $B-1$.
Random multiplicities and spectrum assignments can also change the collision-level
covariance from independent Poisson counts; $z$ is then only a count diagnostic.
See \href{https://docs.scipy.org/doc/scipy/reference/generated/scipy.stats.chisquare.html}
{the Pearson test's sampling assumptions}.
Energy bins use the sample extrema, so unsampled tails are excluded and the
reported inverted yield is an integral over that observed window. A formal test
needs independently specified bins, explicit tail categories, and a covariance or
resampling model appropriate to independent collisions. Zero expected density
with observed emissions must be reported rather than discarded. Doubling quadrature
panels gives a convergence indicator, not a rigorous error bound. The stored
figures and tables have not been regenerated for these qualifications.}
 after \documentclass).
\usepackage[margin=1in]{geometry}
\usepackage{amsmath,amssymb,bm}
\usepackage{graphicx}
\usepackage{booktabs}
\usepackage{hyperref}

\newcommand{\dd}{\,\mathrm{d}}
\newcommand{\Ein}{E}
\newcommand{\Eout}{E'}
\newcommand{\Ecm}{E_{c}}
\newcommand{\mucm}{\mu_{c}}
\newcommand{\mulab}{\mu_0}

\newcommand{\acespec}[1]{J.~L.~Conlin and P.~Romano,
\href{https://mcnp-green.lanl.gov/pdf_files/TechReport_2019_LANL_LA-UR-19-29016_ConlinRomano.pdf}
{\emph{A Compact ENDF (ACE) Format Specification}}, LA-UR-19-29016 (2019),
Sec.~#1}
\newcommand{\ENDFspec}[1]{\href{https://www.nndc.bnl.gov/endfdocs/ENDF-102-2023.pdf}
{\emph{ENDF-6 Formats Manual}} (2023), #1}

% Common closing section: how the verification figures are produced
\newcommand{\verificationmethod}{%
Each figure is produced by \texttt{verification/verify\_laws.py}. For an incident
energy $\Ein$ along $+z$, $10^6$ collisions are sampled with MC/DC's own reaction
sampler (\texttt{sample\_elastic\_scattering}, \texttt{sample\_inelastic\_scattering}
or \texttt{sample\_fission}), and every emitted neutron is histogrammed in
$48 \times 24$ bins of $(\Eout, \mulab)$, with $\mulab$ the lab cosine to $+z$. The
inverted kernel used by RMC is integrated over the same bins (Gauss--Legendre panels;
two-body lines are split where $\mulab^*(\Eout)$ crosses a cosine edge), multiplied by
the number of collisions, and compared with a Pearson-type $\chi^2$ diagnostic
over bins with at least 5 expected counts (the rest pooled into one bin).
The existing tables list the script's $\chi^2/\mathrm{dof}$ and nominal $p$-value,
the emitted neutrons per collision (sampled vs.\ the integral of the
inverted kernel), and the largest relative change of a bin integral when the
quadrature panels are doubled. Left panels: outgoing-energy marginal; middle:
lab-cosine marginal; right: per-bin $z = (\text{sampled} - \text{inverted}) /
\sqrt{\text{inverted}}$.

\noindent These figures check agreement with the MC/DC sampler, not independent
physical validity. Their existing $p$-values are not calibrated goodness-of-fit
probabilities: the script uses $B$ degrees of freedom, whereas a fixed-total
multinomial histogram with $B$ categories and no fitted parameters uses $B-1$.
Random multiplicities and spectrum assignments can also change the collision-level
covariance from independent Poisson counts; $z$ is then only a count diagnostic.
See \href{https://docs.scipy.org/doc/scipy/reference/generated/scipy.stats.chisquare.html}
{the Pearson test's sampling assumptions}.
Energy bins use the sample extrema, so unsampled tails are excluded and the
reported inverted yield is an integral over that observed window. A formal test
needs independently specified bins, explicit tail categories, and a covariance or
resampling model appropriate to independent collisions. Zero expected density
with observed emissions must be reported rather than discarded. Doubling quadrature
panels gives a convergence indicator, not a rigorous error bound. The stored
figures and tables have not been regenerated for these qualifications.}

\usepackage{amsthm}

\newtheorem{proposition}{Proposition}
\newcommand{\tpsi}{\tilde\psi}
\newcommand{\teps}{\tilde\epsilon}
\newcommand{\Sbar}{\bar S}
\newcommand{\Sigt}{\Sigma_t}
\newcommand{\E}{\mathbb{E}}

\title{Sampling the full continuous-energy residual without a two-phase schedule\\
\large analysis, proposal and a 0D prototype}
\date{}

\begin{document}
\maketitle

\noindent Notation as in \texttt{rmc\_math.tex}. $E'$ is the incident and $E$ the
outgoing energy; $T=\mathcal S\tpsi$ is the exact continuous-energy in-scatter of the
trial function. The prototype is \texttt{examples/rmc/residual\_sampling\_0d}; every
number below is reproduced by its scripts (\texttt{results\_*.txt}).

\section{What stalls, and why a better sampler alone cannot fix it}

Each RMC pass estimates $\teps=\Pi L^{-1}r(\tpsi)$ with $\Pi$ the projection onto the
trial space. For any unbiased signed-weight estimator the iterate noise is bounded below
by the noise of the response to the residual; for a source proposal $q$ the source-weight
second moment is $\int r^2/q\ge\|r\|_1^2$ (\texttt{rmc\_math.tex}, eq.~(variance)).
With the full residual, $\E[\tpsi^{(n+1)}\mid\tpsi^{(n)}]=\Pi\psi$, so the iterates
fluctuate around $\Pi\psi$ with a noise level set by
\begin{equation}
  \|r(\Pi\psi)\|_1 = \|q+\mathcal S\Pi\psi-\Sigt\Pi\psi-\mu\partial_z\Pi\psi\|_1 ,
\end{equation}
the residual of the \emph{best} trial function. It is zero only if $\psi$ lies in the
trial space. The stall is that floor: a sampler can approach $\|r\|_1^2$ (the integrated
sampler does), but no sampler goes below it. Two terms keep it large for the
piecewise-constant/linear \emph{flux} basis in continuous energy:
\begin{enumerate}
  \item the in-bin shape of $\Sigt(E)$ in $\Sigt(E)\tpsi$ (zero bin integral, never
  zero pointwise), and
  \item the in-bin shape of the in-scatter, $r_s=T-\Sbar$.
\end{enumerate}
The two-phase schedule works because phase~1 drops both terms (binned $\bar\Sigma_t$,
binned $\Sbar$) and converges to a residual that can vanish. That fixed point carries
the flat-weighting bias, and phase~2 removes it at the floor's noise, averaged as
$1/\sqrt{\text{passes}}$. So the goal is not a cleverer sampler for $r_s$. It is
\textbf{a trial space in which the full residual, scattering integral included, becomes
small}, plus a sampler that draws that residual directly at $\|r\|_1$ cost.

\paragraph{The flat fixed point can be badly wrong.} With a narrow absorption resonance,
the flat-weighted $\bar\Sigma_a$ of the resonance bin is the infinite-dilution value, so
phase~1 converges to a solution with far too much absorption. The prototype problem
(Sec.~\ref{sec:proto}) has a 46\% rms error at the flat $P_0$ fixed point. The phase-1
error operator $I-\Pi L^{-1}L_h$ is also nearly non-contracting there: $\|\teps\|$
oscillates and falls only $\sim 15\%$ per iteration. Phase~2 then has to remove a
large bias with a large-floor sampler.

\section{Proposal}

\subsection{Collision-density trial space}
Use the Petrov--Galerkin space of \texttt{self\_shielded\_basis.tex}: on bin $g$,
\begin{equation}
  \tpsi(z,E,\mu)=\frac{1}{\Sigt(z,E)}\sum_{b=0}^{p}c_{kgjb}\,P_b(x_g(E)),
  \qquad\text{test functions }P_a(x_g(E)) .
\end{equation}
Then $\Sigt\tpsi$ is a polynomial on every bin, and in a cell of a homogeneous material
\begin{equation}
  r = q + T - \textstyle\sum_b c_b P_b \;(-\,\mu\partial_z\tpsi).
\end{equation}
Term 1 above is gone identically, with pointwise $\Sigt$ and no binned removal. The
persistent residual is $(I-\Pi_p)(q+T)$, the degree-$p$ projection error of the emission
density, which is smooth wherever the narrow-resonance picture holds. The flux keeps all
of its $1/\Sigt$ resonance structure. Updates use the existing \texttt{flux} and
\texttt{flux-energy-slope} tallies: $t_a=\int_g\teps P_a$, then $c\leftarrow c+B_g^{-1}t$
with $B_{g,ab}=\int_g P_aP_b/\Sigt$. ($p=2$ needs one new score, $P_2(x_g(E))$ times the
track length.)

\subsection{Unsplit pointwise residual from a tabulated in-scatter kernel}
Since $\tpsi$ is a fixed linear combination of basis functions $\varphi_{\beta'}$,
\begin{equation}
  T(z,E,\mu)=\sum_{\beta'}c_{\beta'}\,U_{\beta'}(E,\mu),\qquad
  U_{\beta'}(E,\mu)=\int_{g'}\dd E'\,\varphi_{\beta'}(E')\sum_rN_r\sigma_r(E')\,
  \bar\tau_{r,j'}(E'\!\to E,\mu),
\end{equation}
and $U_{\beta'}$ does not depend on $\tpsi$. Tabulate it \emph{once} per material, at
outgoing nodes inside each bin. Use piecewise-polynomial interpolation between that
incident bin's own breakpoints: the kinematic edges $\alpha E'_{g'\pm1/2}$,
$E'_{g'\pm1/2}$, thresholds and the kinks of $\Pi_{j'}$ in $\mu$. These are the same
kernel evaluations \texttt{in\_scatter\_density} performs per sample today, done at a
few thousand nodes instead of at every sample of every pass. Kinematic pruning makes
the table banded (elastic) or low-rank (fission: $\chi(E)$ times an incident integral).
Each pass then:
\begin{enumerate}
  \item evaluates $T=\sum c\,U$ (sparse mat-vec) and the \emph{full} residual
  $r=q+T-\Sigt\tpsi$ at the nodes;
  \item builds a proposal $q_x\propto|r|$ on node intervals, mixed with a defensive
  volume-proportional floor (2\% in the prototype), so $q_x>0$ wherever $r\neq0$;
  \item samples $x\sim q_x$ and banks weight $r(x)/(N q_x(x))$, with $r(x)$ evaluated at
  the sampled point (interpolated $T$, exact $\Sigt$ and $\tpsi$).
\end{enumerate}
This removes the $\pm\Sbar$ add-and-subtract, the collision/scatter budget split (audit
finding~4), the per-sample adaptive $E'$ quadrature, and the pilot-floor support
heuristic (audit finding~2). The source-weight second moment is $\approx\|r\|_1^2$ for
the \emph{whole} residual. The table error $\hat U-U$ is a deterministic source error,
so set its tolerance against $\|r\|_1$, not $|T|$ (the absolute-scale point of
\texttt{scattering\_correction\_samplers.tex}, Sec.~4). Moments are optional: computing
$M$ from the same table makes $\int_\beta(\hat T-\Sbar)=0$ exact. $M$ is needed only for
a deterministic Galerkin starting guess.

\subsection{Single loop, residual-driven refinement}
Sample the full residual from the first iteration. Once $\|\teps\|$ stops falling, report
the running mean of the iterates (Polyak averaging; unbiased, since every iterate has
mean $\Pi\psi$). Use the per-bin masses $\int_\beta|r|$, already computed for the
proposal, as an a~posteriori error indicator: split the bins holding the top half of the
mass, project, continue. This is the ECMC h-adaptivity of Peterson, Morel and Ragusa,
applied in energy. In the prototype the mass concentrates where theory says it should:
the first-collision edges $\alpha E$ of the source bin (Placzek transients) and the
absorption resonance with its shadow just below. Kinematic edges of known sources can be
made bin edges from the start.

\section{0D prototype}\label{sec:proto}

Infinite medium, classical elastic scattering isotropic in the COM frame, $A=12$, energy
window $[1,1000]$; $\Sigma_s=1$ plus a scattering resonance (peak 30, half-width 1.5 at
300); $\Sigma_a=0.02\sqrt{100/E}$ plus an absorption resonance (peak 15, half-width 0.5 at
50); source uniform on the top bin $[794.3,1000]$; 30 logarithmic bins. Reference: the
collision-density equation marched on $\sim$18k nodes (grid-converged to $2\times10^{-6}$;
SMC agrees to its noise, $5\times10^{-4}$ at $2\times10^6$ histories). Transport is the
same expected-value 0D random walk for all methods. Budget: 40 iterations $\times$ 20{,}000
histories, the same as SMC with $8\times10^5$. Phase~2 and Polyak averaging start after
iteration 10. Mean of 3 independent repeats:

\begin{center}
\begin{tabular}{llr r}
\toprule
trial space & residual / schedule & rms error & $\|r\|_1$ at end \\
\midrule
--- & SMC, same histories & $9.6\times10^{-4}$ & --- \\
flat $P_0$ & phase 1 only (fixed point) & $4.6\times10^{-1}$ & 3.77 \\
flat $P_0$ & split, every iteration & $1.1\times10^{-2}$ & 5.85 \\
flat $P_0$ & unsplit, every iteration & $7.3\times10^{-3}$ & 5.53 \\
flat $P_0$ & phase 1 + phase 2 (split) & $4.3\times10^{-3}$ & 3.72 \\
collision $P_1$ & phase 1 only (fixed point) & $3.1\times10^{-3}$ & 0.166 \\
collision $P_1$ & split, every iteration & $2.5\times10^{-4}$ & 0.168 \\
collision $P_1$ & unsplit, every iteration & $4.2\times10^{-4}$ & 0.166 \\
collision $P_1$ & phase 1 + phase 2 (unsplit) & $4.0\times10^{-4}$ & 0.166 \\
collision $P_2$ & split, every iteration & $1.2\times10^{-4}$ & 0.055 \\
collision $P_2$ & unsplit, every iteration & $1.5\times10^{-4}$ & 0.054 \\
collision $P_2$ & phase 1 + phase 2 (unsplit) & $1.1\times10^{-4}$ & 0.051 \\
\bottomrule
\end{tabular}
\end{center}

\noindent Residual-driven refinement (collision $P_2$, single loop, unsplit; error on the
original 30 bins):

\begin{center}
\begin{tabular}{rrrr}
\toprule
bins & $\|r\|_1$ at fixed point & fixed-point bias & RMC rms error \\
\midrule
30 & 0.051 & $7.3\times10^{-4}$ & $1.7\times10^{-4}$ \\
33 & 0.025 & $3.3\times10^{-4}$ & $4.5\times10^{-5}$ \\
36 & 0.011 & $6.1\times10^{-5}$ & $1.8\times10^{-5}$ \\
41 & 0.0065 & $2.9\times10^{-5}$ & $1.7\times10^{-5}$ \\
46 & 0.0038 & $1.1\times10^{-5}$ & $1.1\times10^{-5}$ \\
\bottomrule
\end{tabular}
\end{center}

\paragraph{Reading the results.}
\begin{itemize}
  \item The trial space is the lever. Moving from flat $P_0$ to collision-density $P_2$
  cuts the persistent residual 100-fold and the error 30--90-fold at equal cost. The
  phase-1 contraction improves from $\sim0.85$ to $\sim0.1$ per iteration.
  \item Once the basis is right, the schedule hardly matters: single-loop full-residual
  iterations with averaging match phase~1 + phase~2. The phases become optional.
  \item Split vs.\ one-term sampling is measured directly in
  Sec.~\ref{sec:split}: one term is never worse, and the current fixed 90/10 history
  split costs up to $10\times$ in variance.
  \item With 16 refined bins, RMC reaches $1.1\times10^{-5}$ against SMC's
  $9.6\times10^{-4}$ for the same histories: about 90$\times$ lower error, or roughly
  $10^4\times$ in figure of merit at equal histories, before setup cost (the per-history
  transport cost of RMC and SMC was not compared).
\end{itemize}


\section{Split vs.\ one-term sampling: per-pass variance}\label{sec:split}

\texttt{split\_vs\_unsplit.py} and \texttt{split\_vs\_unsplit\_coll.py} run 200
independent passes of $2\times10^4$ histories at a \emph{fixed} trial function (no
feedback between passes, as in a phase-2 correction pass). They report the summed per-bin
variance of the flux correction divided by that of SMC with the same histories. Three
samplers of the same volumetric residual are compared:
\begin{itemize}
  \item \emph{split, 90/10}: $r_c$ and $r_s$ sampled separately with
  \texttt{collision\_fraction}\,=\,0.9 (the \texttt{feature/RMC} default);
  \item \emph{split, mass-proportional}: the same, with histories $\propto\|r_c\|_1$,
  $\|r_s\|_1$ (the best allocation of a split);
  \item \emph{one term}: $r=q+T-\Sigt\tpsi$ sampled $\propto|r|$.
\end{itemize}
All three sample each term $\propto$ its own $|r|$ with the same defensive floor, and
their means agree (max $|\Delta|$/s.e.\ over 30 bins: 1.3--2.3).

\begin{center}
\begin{tabular}{lll rrr}
\toprule
problem & basis & state & 90/10 & mass-prop. & one term \\
\midrule
resonance & flat $P_0$ & early iterate & 116 & 63 & 45 \\
resonance & flat $P_0$ & phase-1 fixed point & 86 & 49 & 33 \\
resonance & flat $P_1$ & early iterate & 31 & 35 & 24 \\
smooth & flat $P_0$ & phase-1 fixed point & 143 & 14.2 & 14.8 \\
smooth & flat $P_1$ & phase-1 fixed point & 0.97 & 0.10 & 0.085 \\
resonance & collision $P_1$ & fixed point & 1.19 & 0.099 & 0.101 \\
resonance & collision $P_2$ & fixed point & 0.138 & 0.0147 & 0.0148 \\
smooth & collision $P_2$ & fixed point & 0.022 & 0.0025 & 0.0025 \\
\bottomrule
\end{tabular}
\end{center}
\noindent (Variance relative to SMC per pass; $<1$ means the pass is less noisy than SMC.
Flat $P_1$ phase~1 diverges on the resonance problem, so its fixed point is not listed.)

\paragraph{Why.} With proportional allocation the split's best source second moment is
$(\|r_c\|_1+\|r_s\|_1)^2$; one term gives $\|r_c+r_s\|_1^2$, never larger by the triangle
inequality. The gain is the cancellation between $r_c$ and $r_s$ inside a bin. With the
flat basis on the resonance problem, $\Sigt(E)\tpsi$ in $r_c$ partly cancels the in-scatter
shape in $r_s$ ($\|r\|_1=3.76$ against $2.55+1.65=4.20$), worth 30\% in variance. Where
there is no cancellation ($r_c\approx0$: smooth cross sections, or the collision basis at
its fixed point) the two coincide. The fixed split is a separate and larger loss: near
convergence $r_s$ holds 95--100\% of the residual mass but gets 10\% of the histories, so
its weights are $\sim10\times$ too large.

\paragraph{For 3D.} Sampling the volumetric residual as one term costs nothing and gives
$1.5$--$10\times$ lower variance than the current scheme. Face terms stay separate
(surface sources), with histories proportional to their mass. The larger lever is the
residual magnitude itself: per pass, flat $P_0$ is 15--90$\times$ \emph{noisier} than
SMC, collision-density $P_2$ is 70--400$\times$ \emph{quieter}, with either
mass-proportional sampler.


\section{Higher energy order: accuracy, cost and figure of merit}\label{sec:order}

\texttt{order\_sweep.py} solves the Galerkin system directly for each order (no Monte
Carlo) and then times one-term residual passes and single-loop runs. FOM
$=1/(\text{error}^2\times\text{time})$ relative to SMC with the same histories; time is
measured wall time (residual build, sampling, transport; transport dominates). Resonance
problem, 30 bins, $2\times10^4$ histories per pass, 40 iterations for the loop:

\begin{center}
\begin{tabular}{l rrr rrr r}
\toprule
 & FP bias & $\|r\|_1$ & pass var & pass time & pass FOM & loop FOM & loop error\\
\midrule
flat $P_0$ & $4.6\times10^{-1}$ & 3.76 & 40 & 1.05 & 0.02 & 0.004 & $1.2\times10^{-2}$\\
flat $P_1$ & $1.5\times10^{-1}$ & 1.45 & 6.4 & 1.12 & 0.14 & 0.03 & $4.4\times10^{-3}$\\
flat $P_2$ & $4.0\times10^{-2}$ & 1.00 & 3.1 & 1.39 & 0.23 & 0.03 & $3.6\times10^{-3}$\\
collision $P_0$ & $3.7\times10^{-2}$ & 1.42 & 9.8 & 1.10 & 0.09 & 0.08 & $2.4\times10^{-3}$\\
collision $P_1$ & $3.1\times10^{-3}$ & 0.166 & 0.13 & 1.29 & 6.0 & 2.8 & $4.0\times10^{-4}$\\
collision $P_2$ & $7.3\times10^{-4}$ & 0.051 & 0.020 & 1.22 & 40 & 34 & $1.1\times10^{-4}$\\
collision $P_3$ & $1.5\times10^{-4}$ & 0.027 & 0.0048 & 1.59 & 131 & 151 & $5.1\times10^{-5}$\\
collision $P_4$ & $1.4\times10^{-5}$ & 0.014 & 0.0015 & 1.80 & 372 & 238 & $3.9\times10^{-5}$\\
\bottomrule
\end{tabular}
\end{center}
\noindent (pass var and pass time relative to an SMC pass; SMC with $8\times10^5$
histories: $7.1\times10^{-4}$ rms.) On the smooth problem the two bases coincide, and the
pass FOM is 7.8, 191, 507, 1594 for $P_1$..$P_4$ (\texttt{results\_order\_smooth.txt}).

\paragraph{Cost.} Going from $P_1$ to $P_2$ costs one more tally score per track and one
more basis evaluation per source point. The time per pass changed by 0.95--1.25$\times$
here, while the variance fell 6--30$\times$. The moments and the in-scatter table use the
same kernel evaluations, weighted by $P_{a'}(x_{g'}(E'))$ and $P_a(x_g(E))$, so setup
grows only through the $(p+1)^2$ storage. The FOM gain from $P_1$ to $P_2$ (7--25$\times$)
tolerates a large per-score tally cost in MC/DC before it disappears.

\paragraph{Stability at low histories per bin} (\texttt{low\_histories.py}). The update
is noisier for higher moments (the $P_a$ tally carries a factor $2a+1$), so the
iteration's noise gain grows with $p$. All orders were stable at $\ge 33$ histories per
energy bin per iteration. At 10 per bin, $P_2$ lost accuracy, $P_3$ diverged in 3 of 8
smooth runs, and $P_4$ in 8 of 8. The observed threshold is roughly 5 histories per
coefficient per iteration. In 3D the count that matters is histories per
phase-space bin (cell $\times$ energy $\times$ angle), so this, not tally cost, limits
the order.

\paragraph{Fewer bins at higher order} (\texttt{bins\_vs\_order.py}, 30 coefficients per
cell). Pointwise flux $L^1$ error, resonance problem: 30$\times P_0$ 6.4\%,
15$\times P_1$ 2.7\%, 10$\times P_2$ 2.4\%, 6$\times P_4$ 2.6\%. Smooth problem:
5.8\%, 2.0\%, 1.4\%, 1.3\%. Order buys accuracy per unknown up to $P_2$. Beyond that,
narrow resonances need bins rather than order.

\paragraph{Recommendation.} Collision-density $P_2$: 7--25$\times$ the FOM of $P_1$ for
about the same time per pass, and stable down to $\sim$30 histories per bin. $P_3$ and
$P_4$ gain another 3--10$\times$ when histories per bin are plentiful, but are the first
to diverge when they are not. Use them only where the phase-space bin count leaves
$\gtrsim 50$ histories per bin. The flat basis needs $P_4$ to reach SMC's FOM on the
resonance problem; higher order does not replace the collision-density weighting.

\section{Caveats and what is untested}
\begin{itemize}
  \item 0D, synthetic data, isotropic COM elastic only. Free-gas upscatter, anisotropic
  and continuum laws and fission change only the table $U$, not the algorithm, but are
  untested here.
  \item In 1D the streaming term of a $1/\Sigt$-weighted basis gives face jumps
  $D(E)=F_R(E)/\Sigma_{t,R}(E)-F_L(E)/\Sigma_{t,L}(E)$ with energy structure at material
  interfaces. They are sampleable exactly (closed form in $E$), but their size relative
  to the flat-basis face residual is unmeasured. The 1D plateau in \texttt{NOTES.md}
  Secs.~2b and~9 is a separate, spatial/angular, issue.
  \item Timings are from the prototype (numba transport, numpy residual); the per-score tally
  cost in MC/DC 3D mesh tallies is unmeasured.
  \item $\Sigt>0$ is required (no vacuum cells in this basis; use the flux basis there).
  $U$ and $M$ are per material, not per nuclide.
  \item The table $\hat U$ is a deterministic approximation; refinement can only drive the
  error down to the table accuracy.
\end{itemize}

\section{Implementation sketch in \texttt{mcdc/rmc}}
\begin{enumerate}
  \item \texttt{kernel.py}/\texttt{driver.py}: build $U_{\beta'}$ at outgoing nodes per
  material with weight $\varphi_{\beta'}=P_{b}/\Sigt$, reusing the inner evaluators of
  \texttt{in\_scatter\_density}; cache it next to the moments (same provenance key plus
  the weight and node set).
  \item \texttt{residual.py}: full residual at nodes; bin masses; node-interval proposal
  with a defensive floor.
  \item \texttt{source.py}: \texttt{sample\_residual} (one term) replacing
  \texttt{sample\_collision\_edge} + \texttt{sample\_correction\_integrated} for that
  basis; face terms as today.
  \item \texttt{driver.py}: Petrov--Galerkin update $B_g^{-1}t$ in
  \texttt{\_read\_epsilon}; single-loop schedule with plateau detection and iterate
  averaging; optional refinement rounds driven by the residual masses.
\end{enumerate}

\begin{thebibliography}{9}
\bibitem{peterson} J.~R.~Peterson, J.~E.~Morel and J.~C.~Ragusa, \emph{Residual Monte
  Carlo for the One-Dimensional Particle Transport Equation}, SIAM J.\ Sci.\ Comput.\
  \textbf{38}(6), B941 (2016); J.~R.~Peterson, \emph{Exponentially-convergent Monte Carlo
  for the One-dimensional Transport Equation}, thesis, Texas A\&M
  (space--angle h-adaptivity after trial-space stagnation).
\bibitem{larsen} M.~Larsen, J.~Bolding, T.~Palmer and J.~E.~Morel, \emph{A Residual
  Monte Carlo Algorithm for Continuous Energy Neutron Transport with Elastic
  Scattering}, Nucl.\ Sci.\ Eng.\ \textbf{200}(3) (2026).
\bibitem{vermaak} J.~I.~C.~Vermaak and J.~E.~Morel, \emph{Transport Error Estimation
  Using Residual Monte Carlo}, J.\ Comput.\ Phys.\ (2022).
\bibitem{hesterberg} T.~Hesterberg, \emph{Weighted average importance sampling and
  defensive mixture distributions}, Technometrics \textbf{37}, 185 (1995).
\end{thebibliography}

\end{document}
